% ch10_f §10.12 (书页 372–374 / p386–p388)
%--E-TAIL--
%JOIN%
必然后果：Heisenberg 哈密顿量 (10.29) 在磁场方向的转动下具有不变性。

从宏观上看，这一点是显而易见的：在长波长的自旋波中，介质在局域上被铁磁性地极化，相邻点之间相对自旋方向不会发生突然的反转。因此，磁振子是一种集体激发，其中电子的自旋在空间中强烈地相互关联。由于这个原因，磁振子–中子相互作用的一般理论并不依赖于电子是局域的还是巡游的（itinerant），因为衍射现象只依赖于这类双粒子的空间与时间关联（§2.8）。反铁磁性的自旋密度波（spin-density wave）理论（§10.7）中隐含的自旋关联，同样能为分析晶体中的磁有序图样给出正确的中子衍射结果（§2.6）。

\section{反铁磁基态}

初看之下，把 §10.11 的方法应用于 Heisenberg 反铁磁体似乎并不困难。假设我们有一个由 Ising 自旋构成的格子，它具有由变量 $\sigma_l$ 所定义的有序排列：在“自旋向上”的格点上 $\sigma_l$ 取 $+1$，在“自旋向下”的格点上取 $-1$。于是我们几乎可以完全仿照 (10.105) 写出我们的哈密顿量
\begin{equation*}
\mathcal{H} = -\sum_{ll'} J_{ll'} \bigl\{ \mathsf{S}_l^z \mathsf{S}_{l'}^z + \tfrac{1}{2}(\mathsf{S}_l^+ \mathsf{S}_{l'}^- + \mathsf{S}_l^- \mathsf{S}_{l'}^+) \bigr\} - \beta \sum_l (\sigma_l H_A \mathsf{S}_l^z + \mathbf{H}\cdot\mathsf{S}_l). \tag{10.124}
\end{equation*}

在这个表达式中，除了外磁场 $\mathbf{H}$ 之外，我们还纳入了一个“各向异性场”（anisotropy field）$H_A$，它总是平行于局域有序自旋的方向。于是，$H_A$ 在“向上”的格点上“向上”，在“向下”的格点上“向下”。这个场给出了量子化 $z$ 轴的方向（它不一定就是外加场的方向），除此之外它是任意的；我们很快就会看到需要它的理由。
%FIG{196}: 反铁磁系统的量子化轴。

现在我们相对于对有序“Ising”态的偏离来定义局域的湮没算符与产生算符。在 $\sigma_l$ 为 $+1$ 的格点上，定义与 (10.109) 相同；而在“自旋向下”的格点上，我们要互换 $\mathsf{S}_l^+$ 与 $\mathsf{S}_l^-$ 的角色，因为我们是从第 $l$ 个离子的态 $|{-S}\rangle_l$ 出发来计量的。于是，在每一个格点上——无论“向上”还是“向下”——如下方程
\begin{align*}
\left.
\begin{aligned}
\mathsf{S}_l^x &\approx (\tfrac{1}{2}S)^{\frac{1}{2}}(a_l + a_l^*), \\
\mathsf{S}_l^y &\approx -i\sigma_l (\tfrac{1}{2}S)^{\frac{1}{2}}(a_l - a_l^*), \\
\mathsf{S}_l^z &= \sigma_l (S - a_l^* a_l),
\end{aligned}
\right\} \tag{10.125}
\end{align*}
定义了湮没和产生自旋偏离的算符 $a_l$ 与 $a_l^*$。（在某些论述中，人们在两个不同的子晶格（sublattice）上定义不同的算符；不过由于当 $l \neq l'$ 时 $a_l$ 与 $a_l^*$ 对易，这只是标签的问题。）

把 (10.125) 代入 (10.124) 是相当容易的。第一项就是这个态的“Ising”能量
\begin{equation*}
\mathcal{E}_0 = -\sum_{ll'} J_{ll'}\, \sigma_l \sigma_{l'} S^2 - N\beta S H_A \tag{10.126}
\end{equation*}
这一项我们可以丢弃。其余的项——在 $a_l$、$a_l^*$ 中保留到二阶乘积——虽然复杂，但很容易写出。

仿照 (10.114) 的一个 Fourier 变换引入新的湮没与产生算符 $b_{\mathbf{q}}$ 和 $b_{\mathbf{q}}^*$，使得
\begin{equation*}
a_l = N^{-\frac{1}{2}} \sum_{\mathbf{q}} b_{\mathbf{q}}\, e^{i\mathbf{q}\cdot\mathbf{l}}, \qquad a_l^* = N^{-\frac{1}{2}} \sum_{\mathbf{q}} b_{\mathbf{q}}^*\, e^{-i\mathbf{q}\cdot\mathbf{l}}. \tag{10.127}
\end{equation*}
除了“Ising”能量 (10.126) 之外，我们发现哈密顿量可以写成
\begin{equation*}
\mathcal{H} = \tfrac{1}{2} \sum_{\mathbf{q}} \bigl\{ A_{\mathbf{q}} (b_{\mathbf{q}}^* b_{\mathbf{q}} + b_{-\mathbf{q}}^* b_{-\mathbf{q}}) + B_{\mathbf{q}} (b_{\mathbf{q}} b_{-\mathbf{q}} + b_{\mathbf{q}}^* b_{-\mathbf{q}}^*) \bigr\} + \mathbf{H}\cdot\mathbf{M}, \tag{10.128}
\end{equation*}
其中系数与 (2.12) 或 (10.116) 相类似：——
\begin{align*}
\left.
\begin{aligned}
A_{\mathbf{q}} &= \beta H_A + \sum_{\mathbf{h}} 2S J(\mathbf{h}) \bigl\{ \sigma_{\mathbf{h}} - \tfrac{1}{2}(1 + \sigma_{\mathbf{h}})\cos\mathbf{q}\cdot\mathbf{h} \bigr\}, \\
B_{\mathbf{q}} &= -\sum_{\mathbf{h}} S J(\mathbf{h})(1 - \sigma_{\mathbf{h}})\cos\mathbf{q}\cdot\mathbf{h}.
\end{aligned}
\right\} \tag{10.129}
\end{align*}
在这些表达式中，$\sigma_{\mathbf{h}}$ 表示位于 $\mathbf{h}$ 处的格点相对于中心格点符号的符号。我们可以把 $J(\mathbf{h})$ 的力程比最近邻相互作用更长、甚至在某些距离上为铁磁型的情形也包括进来；唯一必要的条件是 $\sigma_l$ 应当正确地描述 Ising 模型的基态。

在 (10.128) 中，外磁场与一个复杂的算符 $\mathbf{M}$ 联结在一起出现；我们不打算一般地定义它，尽管在简单的情形中可以看出，它对应于一个在两个子晶格的格点上（沿 $z$ 轴）看似正负交替的矩。

这是我们在 (10.125) 中对局域量子化轴的任意选取所造成的结果——对这个项我们就不再深究了。

我们的“交替”态的另一个效应是：(10.128) 中存在把 $b_{\mathbf{q}}$ 与 $b_{-\mathbf{q}}$ 联结起来的项。当我们把 (10.125) 代入 (10.124) 时，它们通过像 $a_l a_{l'}$ 这样的乘积产生，而这类乘积在铁磁态 (10.113) 中并不出现。强加上一个 Ising 自旋背景系统降低了体系的基本平移对称性，因此哈密顿量不再能被 Fourier 变换 (10.127) 自动对角化。或许我们应当取每个晶胞两个自旋，并作一种颇不相同的 Fourier 分析。

然而，这个问题实际上并不比晶格动力学的类似问题中相互作用矩阵 (2.13) 的约化更困难。我们通过一个正则变换引入新变量 $c_{\mathbf{q}}$、$c_{\mathbf{q}}^*$：
\begin{align*}
\left.
\begin{aligned}
b_{\mathbf{q}} &= c_{\mathbf{q}}\cosh\theta_{\mathbf{q}} + c_{-\mathbf{q}}^*\sinh\theta_{\mathbf{q}}, \qquad b_{\mathbf{q}}^* = c_{\mathbf{q}}^*\cosh\theta_{\mathbf{q}} + c_{-\mathbf{q}}\sinh\theta_{\mathbf{q}}, \\
b_{-\mathbf{q}} &= c_{\mathbf{q}}^*\sinh\theta_{\mathbf{q}} + c_{-\mathbf{q}}\cosh\theta_{\mathbf{q}}, \qquad b_{-\mathbf{q}}^* = c_{\mathbf{q}}\sinh\theta_{\mathbf{q}} + c_{-\mathbf{q}}^*\cosh\theta_{\mathbf{q}},
\end{aligned}
\right\} \tag{10.130}
\end{align*}
它保持对易子
\begin{equation*}
[c_{\mathbf{q}}, c_{\mathbf{q}}^*] = \cosh^2\theta_{\mathbf{q}} - \sinh^2\theta_{\mathbf{q}} = 1. \tag{10.131}
\end{equation*}
参数 $\theta_{\mathbf{q}}$ 表示在这四个算符的空间中绕一个虚角度的转动。值得注意的是，$c_{\mathbf{q}}$ 的共轭是 $c_{-\mathbf{q}}^*$；我们必须反转运动方向，才能使它看上去像是给定波的物理共轭。

如果把 (10.130) 代入 (10.128)，就会得到一组“非对角”项
\begin{equation*}
\tfrac{1}{2} \sum_{\mathbf{q}} \bigl\{ 2A_{\mathbf{q}}\cosh\theta_{\mathbf{q}}\sinh\theta_{\mathbf{q}} + B_{\mathbf{q}}(\cosh^2\theta_{\mathbf{q}} + \sinh^2\theta_{\mathbf{q}}) \bigr\} (c_{\mathbf{q}} c_{-\mathbf{q}} + c_{\mathbf{q}}^* c_{-\mathbf{q}}^*), \tag{10.132}
\end{equation*}
只要把 $\theta_{\mathbf{q}}$ 选得满足
\begin{align*}
\frac{B_{\mathbf{q}}}{A_{\mathbf{q}}} &= -\frac{2\cosh\theta_{\mathbf{q}}\sinh\theta_{\mathbf{q}}}{\cosh^2\theta_{\mathbf{q}} + \sinh^2\theta_{\mathbf{q}}} \\
&= -\tanh 2\theta_{\mathbf{q}}, \tag{10.133}
\end{align*}
这组项便可以对每一个 $\mathbf{q}$ 值都变为零。对角项于是给出如下结果
\begin{equation*}
\mathcal{H} = \sum_{\mathbf{q}} \bigl[ (A_{\mathbf{q}}^2 - B_{\mathbf{q}}^2)^{\frac{1}{2}} (c_{\mathbf{q}}^* c_{\mathbf{q}} + \tfrac{1}{2}) - A_{\mathbf{q}} \bigr], \tag{10.134}
\end{equation*}
其中我们已利用对易关系把算符排成了这一顺序。现在的求和遍及 $\mathbf{q}$ 的一切取值，正负兼备。
